\documentclass[11pt]{article}
\usepackage[margin=1in]{geometry}
\usepackage{amsmath,amssymb,amsthm,mathtools}
\usepackage{graphicx}
\usepackage{microtype}
\usepackage[hidelinks]{hyperref}
\usepackage{xcolor}

\newtheorem{theorem}{Theorem}
\newtheorem{lemma}{Lemma}
\newtheorem{corollary}{Corollary}
\newcommand{\Hh}{\mathcal H}
\newcommand{\B}{\mathcal B}
\newcommand{\Sym}{\operatorname{Sym}}
\newcommand{\ess}{\mathrm e}
\newcommand{\od}{\mathbin{\odot}}

\title{Symmetric tensor powers reflect subnormality}
\author{Candidate solution packet for arXiv:1503.03692}
\date{13 August 2026}

\begin{document}
\maketitle

\begin{center}
\fcolorbox{black}{gray!8}{\parbox{0.91\linewidth}{
\textbf{Status.} Candidate full solution to the open question after
Proposition~10.5 of Stochel--Stochel \cite{SS}, pending expert review.  The
operator theorem below also answers the two-fold symmetric-tensor question
left open by Stochel in 1996 \cite{S96}.
}}
\end{center}

\section{Source question}

The source is J.~Stochel and J.~B.~Stochel, \emph{Composition operators on
Hilbert spaces of entire functions with analytic symbols},
arXiv:1503.03692v3 \cite{SS}.  On PDF page~28, after Proposition~10.5, it asks
whether the ``only if'' direction remains true when $\Phi'(0)=0$.  Here is the
source passage, including the proposition to which the question refers.

\begin{center}
\includegraphics[width=0.98\linewidth]{figures/open_problem_crop.png}
\end{center}

In full: for a nonconstant entire function
$\Phi(z)=\sum_{n\ge0}a_nz^n$ with $a_n\ge0$, and $A\in\B(\Hh)$, does
subnormality of $C_A$ (respectively $C_A^*$) imply subnormality of $A^*$
(respectively $A$) without assuming $a_1=\Phi'(0)\ne0$?  For
$\Phi(z)=z^2$, the source identifies this with the 1996 question: does
subnormality of $T^{\od 2}$ imply subnormality of $T$?

\section{Definitions and statement}

Let $P_m$ be the orthogonal projection from $\Hh^{\otimes m}$ onto its
symmetric subspace $\Hh^{\od m}$.  For $T\in\B(\Hh)$, the symmetric power
$T^{\od m}$ is the restriction of $T^{\otimes m}$ to $\Hh^{\od m}$.  A
bounded operator is \emph{subnormal} if it is the restriction of a normal
operator to an invariant closed subspace.

\begin{theorem}[Symmetric-power reflection]\label{thm:main}
Let $T$ be a bounded operator on a complex Hilbert space and let $m\ge1$.
Then
\[
             T^{\od m}\text{ is subnormal}
             \quad\Longleftrightarrow\quad
             T\text{ is subnormal}.
\]
\end{theorem}

The forward implication is new relative to the two source questions.  The
reverse implication is standard: a normal extension of $T$ induces a normal
symmetric-power extension.

\section{Proof intuition}

Pure powers $x^{\otimes m}$ immediately descend hyponormality from
$T^{\od m}$ to $T$.  If this already makes $T$ normal, there is nothing more
to prove.  Otherwise a nonzero boundary point $\lambda$ of the essential
spectrum supplies a weakly-null approximate eigenvector sequence $y_j$.
Insert an arbitrary vector $x$ into one tensor slot and $m-1$ copies of $y_j$
into the others, then symmetrize.  The $k$th orbit norm under $T^{\od m}$ is
a Stieltjes moment sequence in $k$.  As $j\to\infty$, all cross terms vanish
weakly and the remaining expression is
\[
 \frac1m|\lambda|^{2k(m-1)}\lVert T^kx\rVert^2.
\]
Stieltjes moment sequences are closed under pointwise limits, and the harmless
geometric factor can be rescaled away.  Lambert's norm-moment criterion then
forces $T$ itself to be subnormal.

\section{Proof}

We use two standard facts.  First, a bounded operator $R$ is subnormal if and
only if, for every $x$, the sequence
$(\lVert R^kx\rVert^2)_{k\ge0}$ is a Stieltjes moment sequence
\cite[Proposition~2.3]{S96}.  Second, Stieltjes moment sequences are closed
under pointwise convergence \cite[(S2)]{S96}.

\begin{lemma}[A nonzero singular spectral point]\label{lem:singular}
Let $T$ be hyponormal on an infinite-dimensional Hilbert space.  Either $T$
is normal, or there exist $\lambda\ne0$ and unit vectors $y_j\rightharpoonup0$
such that
\[
                     \lVert(T-\lambda)y_j\rVert\longrightarrow0.
\]
\end{lemma}

\begin{proof}
If the essential spectrum is $\{0\}$, then $T$ is a Riesz operator, so its
spectrum is countable.  Putnam's inequality for hyponormal operators gives
\[
 \lVert T^*T-TT^*\rVert
 \le \frac1\pi\operatorname{Area}(\sigma(T))=0,
\]
and $T$ is normal.

Otherwise the compact set $\sigma_{\ess}(T)$ has a nonzero boundary point
$\lambda$.  A boundary point of the Fredholm essential spectrum cannot be an
upper semi-Fredholm point: the upper semi-Fredholm set is open, and a
non-Fredholm upper semi-Fredholm point has a whole neighborhood of
non-Fredholm points (the extended index is locally constant).  Hence
$T-\lambda$ is not upper semi-Fredholm.  The singular-sequence
characterization of upper semi-Fredholm operators now gives unit vectors
$y_j\rightharpoonup0$ with $(T-\lambda)y_j\to0$.
\end{proof}

\begin{lemma}[The symmetrized one-slot identity]\label{lem:norm}
For $u,v\in\Hh$ and $m\ge2$, put
\[
 J_vu=P_m(u\otimes v^{\otimes(m-1)})
     =\frac1m\sum_{r=1}^m
       v^{\otimes(r-1)}\otimes u\otimes v^{\otimes(m-r)}.
\]
Then
\begin{equation}\label{eq:norm}
 \lVert J_vu\rVert^2
 =\frac1m\lVert u\rVert^2\lVert v\rVert^{2m-2}
  +\frac{m-1}{m}|\langle u,v\rangle|^2\lVert v\rVert^{2m-4}.
\end{equation}
Moreover,
\begin{equation}\label{eq:intertwine}
 (T^{\od m})^kJ_vu=J_{T^kv}(T^ku),\qquad k\ge0.
\end{equation}
\end{lemma}

\begin{proof}
In the displayed sum defining $J_vu$, there are $m$ diagonal inner products
and $m(m-1)$ off-diagonal inner products.  Their values are respectively
$\lVert u\rVert^2\lVert v\rVert^{2m-2}$ and
$|\langle u,v\rangle|^2\lVert v\rVert^{2m-4}$, which proves
\eqref{eq:norm}.  Equation \eqref{eq:intertwine} follows because
$T^{\otimes m}$ commutes with every tensor-slot permutation and hence with
$P_m$.
\end{proof}

\begin{proof}[Proof of Theorem~\ref{thm:main}]
Only the forward implication needs proof.  The case $m=1$ is tautological.
Put $S=T^{\od m}$ and suppose $m\ge2$ and $S$ is subnormal.  Then $S$ is
hyponormal.  Applied to the pure symmetric tensor $x^{\otimes m}$, its
hyponormality inequality gives
\[
 \lVert T^*x\rVert^m
 =\lVert S^*x^{\otimes m}\rVert
 \le\lVert Sx^{\otimes m}\rVert
 =\lVert Tx\rVert^m.
\]
Thus $T$ is hyponormal.  If $\Hh$ is finite-dimensional, tracing the
self-commutator shows that every hyponormal operator is normal, and we are
done.  The same is true if $T$ is normal in infinite dimension.

In the remaining case choose $\lambda\ne0$ and $y_j$ from
Lemma~\ref{lem:singular}.  For each fixed $k$,
\begin{equation}\label{eq:powers}
 T^ky_j-\lambda^ky_j\longrightarrow0,
\end{equation}
by the polynomial identity for $T^k-\lambda^k$.  Fix $x\in\Hh$ and set
$\xi_j=J_{y_j}x$.  By subnormality of $S$ and Lambert's criterion,
\[
       a^{(j)}_k:=\lVert S^k\xi_j\rVert^2,qquad k\ge0,
\]
is a Stieltjes moment sequence.  Lemma~\ref{lem:norm} and
\eqref{eq:intertwine} give
\begin{align*}
 a^{(j)}_k
 &=\frac1m\lVert T^kx\rVert^2\lVert T^ky_j\rVert^{2m-2}\\
 &\quad+\frac{m-1}{m}|\langle T^kx,T^ky_j\rangle|^2
                    \lVert T^ky_j\rVert^{2m-4}.
\end{align*}
By \eqref{eq:powers}, $\lVert T^ky_j\rVert\to|\lambda|^k$.  Also
$T^ky_j\rightharpoonup0$, so the second term tends to zero.  Consequently,
for every $k\ge0$,
\begin{equation}\label{eq:limit}
 a^{(j)}_k\longrightarrow
 a_k:=\frac1m|\lambda|^{2k(m-1)}\lVert T^kx\rVert^2.
\end{equation}
The pointwise closure property makes $(a_k)$ a Stieltjes moment sequence.
Multiplication by $m$ and removal of the positive geometric factor
$c^k$, $c=|\lambda|^{2(m-1)}$, correspond respectively to scaling a
representing measure and pushing it forward by $t\mapsto t/c$.  Hence
$(\lVert T^kx\rVert^2)_{k\ge0}$ is a Stieltjes moment sequence.  Since $x$
was arbitrary, Lambert's criterion proves that $T$ is subnormal.
\end{proof}

\section{Application to the source question}

\begin{corollary}\label{cor:composition}
Let $\Phi(z)=\sum_{n\ge0}a_nz^n$ be a nonconstant entire function with
$a_n\ge0$, and let $A\in\B(\Hh)$.  Without any hypothesis on $\Phi'(0)$,
\[
 C_A\text{ is subnormal}\Longleftrightarrow A^*\text{ is subnormal},
 \qquad
 C_A^*\text{ is subnormal}\Longleftrightarrow A\text{ is subnormal}.
\]
\end{corollary}

\begin{proof}
The Fock model of \cite[Theorem~7.2]{SS} identifies $C_A^*$ with the
orthogonal sum of $\Xi_Q(A)^{\od n}$ over $n$ for which $a_n>0$, where
$\Xi_Q(A)=QAQ$ for a conjugation $Q$.  Taking adjoints gives the corresponding
model for $C_A$.  Because $\Phi$ is nonconstant, at least one such index
$n\ge1$ exists.  A reducing summand of a subnormal operator is subnormal, so
Theorem~\ref{thm:main}, and invariance of subnormality under conjugation,
give the two necessity directions.  The sufficiency directions are the
standard symmetric-power construction already recorded in
\cite[Proposition~10.5]{SS}; they never required $\Phi'(0)\ne0$.
\end{proof}

\section{Verification}

The companion script \texttt{code/verify\_symmetrization.py} checked
\eqref{eq:norm} and \eqref{eq:intertwine} in 400 seeded random complex cases
with dimensions $2$--$4$ and powers $m=2$--$4$.  The worst relative error was
below $2\cdot10^{-12}$.  It also used the standard weakly-null approximate
eigenvectors for a unilateral shift at a unimodular $\lambda$ and observed
the residual and the maximum error in \eqref{eq:limit} decrease over lengths
$32,64,\ldots,1024$.  The command was
\begin{verbatim}
conda run --no-capture-output -n sandbox python code/verify_symmetrization.py
\end{verbatim}
and all checks passed.  These finite computations check the tensor algebra
and limiting formula; they do not prove the spectral lemma or subnormality.

\section{Novelty check, limitations, and reviewer focus}

The bounded check was performed on 13 August 2026.  The run indexes were
searched for arXiv:1503.03692, the exact title, ``symmetric tensor power,''
and ``subnormality.''  Web searches used the exact 1996 open-question wording,
the 1996 DOI, the source question with $\Phi'(0)=0$, and combinations of
``subnormal'' with ``symmetric tensor square/power.''  The recent survey
Garcia--O'Loughlin--Yu, arXiv:2306.14807 \cite{GOY}, was also searched; it
contains no occurrence of ``subnormal.''  Results located the original 1996
question and the 2017 source but no claimed solution.  This is a bounded
search, not a proof of novelty.

The theorem concerns bounded operators.  It does not assert an unbounded
symmetric-power reflection theorem.  Corollary~\ref{cor:composition} allows
the composition operator itself to be unbounded/subnormal in the source's
closed-operator sense; the extracted tensor summand is bounded because
$A\in\B(\Hh)$.  A reviewer should check especially (i) the use of the
upper-semi-Fredholm/singular-sequence characterization at a boundary point of
$\sigma_{\ess}(T)$, (ii) the Riesz-hyponormal normality argument when
$\sigma_{\ess}(T)=\{0\}$, and (iii) closure and rescaling of the limiting
Stieltjes moment sequence.

\begin{thebibliography}{9}
\bibitem{SS}
J.~Stochel and J.~B.~Stochel,
\emph{Composition operators on Hilbert spaces of entire functions with
analytic symbols}, Journal of Mathematical Analysis and Applications
\textbf{454} (2017), 1019--1066; arXiv:1503.03692.

\bibitem{S96}
J.~Stochel,
\emph{Seminormality of operators from their tensor product},
Proceedings of the American Mathematical Society \textbf{124} (1996),
135--140.  DOI: \href{https://doi.org/10.1090/S0002-9939-96-03017-1}
{10.1090/S0002-9939-96-03017-1}.

\bibitem{GOY}
S.~R.~Garcia, R.~O'Loughlin, and J.~Yu,
\emph{Symmetric and antisymmetric tensor products for the function-theoretic
operator theorist}, Canadian Journal of Mathematics \textbf{77} (2025),
324--346; arXiv:2306.14807.
\end{thebibliography}

\end{document}
